\documentclass[../main.tex]{subfiles}
\begin{document}
\section{Related Literature} \label{sec:related_lit}
This section describes previous research conducted on sentiment analysis in finance (\ref{sentiment analysis in finance}) and text classification using pre-trained language models (\ref{text classification}).

\subsection{Sentiment analysis in finance} \label{sentiment analysis in finance}
Sentiment analysis is the task of extracting sentiments or opinions of people from written language \cite{Liu2012}. We can divide the recent efforts into two groups: 1) Machine learning methods with features extracted from text with "word counting" \cite{Agarwal2016,Whitelaw2005,conf/icwsm/MartineauF09,Tripathy2016}, 2) Deep learning methods, where text is represented by a sequence of embeddings \cite{Severyn2015,Araque2017,Zhang2018}. The former suffers from inability to represent the semantic information that results from a particular sequence of words, while the latter is often deemed as too "data-hungry" as it learns a much higher number of parameters \cite{2018arXiv180100631M}.  

Financial sentiment analysis differs from general sentiment analysis not only in domain, but also the purpose. The purpose behind financial sentiment analysis is usually guessing how the markets will react with the information presented in the text \cite{Li2014}. Loughran and McDonald (2016) presents a thorough survey of recent works on financial text analysis utilizing machine learning with "bag-of-words" approach or lexicon-based methods \cite{Loughran2016}. For example, in Loughran and McDonald (2011), they create a dictionary of financial terms with assigned values such as "positive" or "uncertain" and  measure the tone of a documents by counting words with a specific dictionary value  \cite{Loughran2011}. Another example is Pagolu et al. (2016), where n-grams from tweets with financial information are fed into supervised machine learning algorithms to detect the sentiment regarding the financial entity mentioned.

On of the first papers that used deep learning methods for textual financial polarity analysis was Kraus and Feuerriegel (2017) \cite{Kraus2017}. They apply an LSTM neural network to ad-hoc company announcements to predict stock-market movements and show that method to be more accurate than traditional machine learning approaches. They find pre-training their model on a larger corpus to improve the result, however their pre-training is done on a labeled dataset, which is a more limiting approach then ours, as we pre-train a language model as an unsupervised task.

There are several other works that employ various types of neural architectures for financial sentiment analysis. Sohangir et al. (2018) \cite{Sohangir2018} apply several generic neural network architectures to a StockTwits dataset, finding CNN as the best performing neural network architecture. Lutz et al. 2018 \cite{Lutz2018} take the approach of using \textit{doc2vec} to generate sentence embeddings in a particular company ad-hoc announcement and utilize multi-instance learning to predict stock market outcomes. Maia et al. (2018) \cite{Maia20182} use a combination of text simplification and LSTM network to classify a set of sentences from financial news according to their sentiment and achieve state-of-the-art results for the Financial PhraseBank, which is used in thesis as well.

Due to lack of large labeled financial datasets, it is difficult to utilize neural networks to their full potential for sentiment analysis. Even when their first (word embedding) layers are initialized with pre-trained values, the rest of the model still needs to learn complex relations with relatively small amount of labeled data. A more promising solution could be initializing almost the entire model with pre-trained values and fine-tuning those values with respect to the classification task.

\subsection{Text classification using pre-trained language models} \label{text classification}

Language modeling is the task of predicting the next word in a given piece of text. One of the most important recent developments in natural language processing is the realization that a model trained for language modeling can be successfully fine-tuned for most down-stream NLP tasks with small modifications. These models are usually trained on very large corpora, and then with addition of suitable task-specific layers fine-tuned on the target dataset \cite{Kant2018}. Text classification, which is the focus of this thesis, is one of the obvious use-cases for this approach.

ELMo (Embeddings from Language Models) \cite{Peters2018} was one of the first successful applications of this approach. With ELMo, a deep bidirectional language model is pre-trained on a large corpus. For each word, hidden states of this model is used to compute a contextualized representation. Using the pre-trained weights of ELMo, contextualized word embeddings can be calculated for any piece of text. Initializing embeddings for down-stream tasks with those were shown to improve performance on most tasks compared to static word embeddings such as word2vec or GloVe. For text classification tasks like SST-5, it achieved state-of-the-art performance when used together with a bi-attentive classification network \cite{McCann2017}. 

Although ELMo makes use of pre-trained language models for contextualizing representations, still the information extracted using a language model is present only in the first layer of any model using it. ULMFit (Universal Language Model Fine-tuning) \cite{Howard2018} was the first paper to achieve true transfer learning for NLP, as using novel techniques such as discriminative fine-tuning, slanted triangular learning rates and gradual unfreezing. They were able to efficiently fine-tune a whole pre-trained language model for text classification. They also introduced further pre-training of the language model on a domain-specific corpus, assuming target task data comes from a different distribution than the general corpus the initial model was trained on. 

ULMFit's main idea of efficiently fine-tuning a pre-trained a language model for down-stream tasks was brought to another level with Bidirectional Encoder Representations from Transformers (BERT) \cite{Devlin2018}, which is also the main focus of this paper. BERT has two important differences from what came before: 1) It defines the task of language modeling as predicting randomly masked tokens in a sequence rather than the next token, in addition to a task of classifying two sentences as following each other or not. 2) It is a very big network trained on an unprecedentedly large corpus. These two factors enabled in to achieve state-of-the-art results in multiple NLP tasks such as, natural language inference or question answering.

The specifics of fine-tuning BERT for text classification has not been researched thoroughly. One such recent work is Sun et al. (2019) \cite{Sun2019}. They conduct a series of experiments regarding different configurations of BERT for text classification. Some of their results will be referenced throughout the rest of the thesis, for the configuration of our model.
\end{document}

